\section{Related works}
\label{sec:related}

This section discusses the work our pipeline builds on and is measured against, grouped by the question each line of work answers. For each work we give the problem, the method and the outcome reported by its authors.

\subsection{Radiance fields and splatting toolkits}
NeRF \cite{mildenhall2020nerf} represents a scene as a neural field rendered by volume rendering; 3DGS \cite{kerbl3dgs} replaces it with millions of anisotropic Gaussians rasterised in real time. Nerfstudio \cite{tancik2023nerfstudio} is a modular framework for such methods, with data importers, a viewer and export tools, and its rasteriser library gsplat \cite{ye2025gsplat} reports up to 10\% less training time and 4$\times$ less memory than the original 3DGS code. 2D Gaussian splatting \cite{huang2024twodgs} improves the poor surfaces of 3DGS by collapsing the Gaussians into oriented disks with depth-distortion and normal-consistency terms, and reports detailed, noise-free geometry. 3DGS-MCMC \cite{kheradmand2024mcmc}, used in our v1, recasts densification as Markov chain Monte Carlo sampling.

\subsection{Gaussian splatting for indoor rooms}
DN-Splatter \cite{turkulainen2025dnsplatter} targets textureless indoor scenes: it supervises the Gaussians with sensor depth through a loss weighted by image gradients, adds a smoothness term and monocular normals, and extracts a mesh directly from the result. AGS-Mesh \cite{ren2025agsmesh} notes that phone depth is coarse and monocular normals are inconsistent across views, filters these priors by their consistency during training and adds a scale-aware meshing step, improving meshes and novel views on indoor room datasets. GaussianRoom \cite{xiang2025gaussianroom} couples a neural signed distance field with 3DGS to guide densification and pruning, with normal and edge priors for textureless areas, and reports state-of-the-art results on ScanNet and ScanNet++. IndoorGS \cite{ruan2025indoorgs} lifts 2D line segments to 3D and derives plane cues for textureless regions to drive initialisation and density control, reporting more accurate geometry and rendering. PlanarGS \cite{jin2025planargs} finds planar regions with a vision-language segmentation model and supervises the Gaussians with planar and depth/normal terms, outperforming earlier methods on indoor benchmarks. LighthouseGS \cite{han2026lighthousegs} handles panorama-style phone captures by using the phone's poses and monocular depth as rough priors, initialising on planes and correcting drift and auto-exposure. PocketGS \cite{guo2026pocketgs} trains 3DGS on the phone itself within minutes. All of these report image quality and surface accuracy; none asks whether a person can walk in the result.

\subsection{Camera poses from images alone}
COLMAP \cite{schoenberger2016colmap} recovers poses by matching local features; learned features such as SuperPoint \cite{detone2018superpoint} with LightGlue \cite{lindenberger2023lightglue} in hloc \cite{sarlin2019hloc}, and visual SLAM such as ORB-SLAM3 \cite{campos2021orbslam3}, rely on the same principle. MASt3R-SfM \cite{duisterhof2025mast3rsfm} builds structure from motion on a 3D foundation model and aligns its local reconstructions at linear cost; VGGT \cite{wang2025vggt} predicts cameras, depth and point maps from up to hundreds of views in one forward pass, in under a second. These newer models depend less on texture than the methods that failed in our v1, but we did not test them; taking poses from the phone removed the problem for this project.

\subsection{Walking and navigating in Gaussian splats}
The SuperSplat Viewer \cite{supersplatviewer} offers a first-person walk mode whose collision is a voxel volume generated from the splat by splat-transform \cite{splattransform, supersplatwalk}. Splat-Nav \cite{chen2025splatnav} plans safe robot trajectories directly in a splat map using collision constraints derived from the Gaussians, tested in 126 hardware flights. NavGSim \cite{liu2026navgsim} is a splat-based simulator for multi-room navigation that extracts navigable areas by slicing the Gaussians. The Neverwhere benchmark \cite{chen2026neverwhere} turns more than sixty splat reconstructions into closed-loop tests for visual legged locomotion, and LucidSim \cite{yu2024lucidsim} trains a visual parkour policy from generated images; in both, physics runs on separate collision geometry while the images come from the splat or a generator. None of these reports how accurately its collision follows the real walls.

\subsection{Capture guidance and view selection}
Kopanas and Drettakis \cite{kopanas2023placement} place cameras for free-viewpoint navigation in rooms using observation frequency and angular uniformity. FisherRF \cite{jiang2024fisherrf} picks the next view by expected information gain from the Fisher information of the model's parameters, while NeRF Director \cite{xiao2024nerfdirector} shows that uniform coverage alone, without any uncertainty estimate, selects good training views. COVER \cite{chen2026cover} shows that an approximation of the Fisher information gain reduces to covering geometry that past cameras saw poorly, and AGILE-GS \cite{khass2026agilegs} cuts next-best-view selection time by one to two orders of magnitude. For people capturing with phones, Zhang et al.\ \cite{zhang2026objectcentered} guide object captures live by spherical coverage and beat free capture with fewer images; Yasunaga et al.\ \cite{yasunaga2025userloop} show in augmented reality where a quick reconstruction is wrong so that the user adds views there; LiMRSF \cite{chang2026limrsf} overlays a LiDAR SLAM point cloud on the real room in a mixed-reality headset to reveal blind spots and double walls while scanning. MUGSQA \cite{chen2026mugsqa} collects human quality ratings of splats reconstructed from inputs of varying number, resolution and viewing distance.

\subsection{Limitations of existing work}
Limitations of the existing methods, as they affect a walkable room, are:
\begin{itemize}
    \item Indoor splatting methods are evaluated on image and mesh quality only; walkability is not tested.
    \item Walk collision built from a splat is the default in a widely used viewer, but its accuracy against a ground-truth scan is not published; robotics work pairs splats with separate collision geometry without measuring the difference.
    \item Pose estimation from images alone still struggles in plain-walled rooms (our v1; the newest feed-forward models were not tested).
    \item Capture feedback is live, object-scale, or needs special hardware or a trained model; none checks an exported room capture before training.
    \item The path from a phone to a walkable room spans separate tools with different coordinate conventions and defaults that do not suit a shared GPU.
\end{itemize}
We address the first, second and last points directly (Sections~\ref{sec:method} and~\ref{sec:collision-results}). The fourth was our first novelty attempt, which failed its own test (Section~\ref{sec:negative}); the third is future work.
